\editionPage[\pdfbookmark[0]{Défis d’approfondissement}{bm-defis-debut-parent}]{defis-debut}{Vingt défis}{Un défi par chapitre · Parcours facultatif}{Chercher avant de consulter les deux indications.}
\begin{blocEdition}{defis-debut}
Ces problèmes demandent davantage d’initiative. Ils utilisent les outils du livre ; les prolongements sont définis dans l’énoncé. Les défis sont facultatifs et approfondissent les outils du programme.\par\bigskip
\begin{savoirfaire}[Une recherche qui avance]Faire un schéma, tester un petit cas, comparer deux méthodes, conserver les essais et préciser ce qui reste à prouver. Ouvrir la première indication seulement après une recherche réelle ; revenir au problème avant de lire la seconde.\end{savoirfaire}
\begin{savoirfaire}[Après la solution]Identifier l’idée décisive et la refaire sur une feuille vierge. Expliquer pourquoi elle fonctionne compte autant que retrouver le résultat.\end{savoirfaire}
\end{blocEdition}
\editionPage{defi-ex-1}{D01 · Défi}{Une commande au meilleur prix}{Approfondissement · 30 à 60 minutes indicatives}
\begin{blocEdition}{defi-ex-1}
On commande 18 carnets. A vend chaque carnet 2,40 euros avec 5 euros de frais par commande. B vend des lots de cinq à 10 euros, sans frais. On peut acheter plus de carnets que nécessaire. Déterminer la commande la moins chère.\par\vspace{5mm}{\otLight\fontsize{8}{10}\selectfont\color{otMuted}Espace de recherche}\par\par\vspace{7mm}{\color{otRule!65}\hrule height .25pt}\par\vspace{7mm}{\color{otRule!65}\hrule height .25pt}\par\vspace{7mm}{\color{otRule!65}\hrule height .25pt}\par\vspace{7mm}{\color{otRule!65}\hrule height .25pt}\par\vspace{7mm}{\color{otRule!65}\hrule height .25pt}\par\vspace{7mm}{\color{otRule!65}\hrule height .25pt}
\end{blocEdition}
\editionPage{defi-ex-2}{D02 · Défi}{Deux restes et une divisibilité}{Approfondissement · 30 à 60 minutes indicatives}
\begin{blocEdition}{defi-ex-2}
Trouver tous les entiers entre 100 et 150 qui laissent un reste de 2 dans la division par 7 et sont divisibles par 9. Expliquer pourquoi la liste est complète.\par\vspace{5mm}{\otLight\fontsize{8}{10}\selectfont\color{otMuted}Espace de recherche}\par\par\vspace{7mm}{\color{otRule!65}\hrule height .25pt}\par\vspace{7mm}{\color{otRule!65}\hrule height .25pt}\par\vspace{7mm}{\color{otRule!65}\hrule height .25pt}\par\vspace{7mm}{\color{otRule!65}\hrule height .25pt}\par\vspace{7mm}{\color{otRule!65}\hrule height .25pt}\par\vspace{7mm}{\color{otRule!65}\hrule height .25pt}
\end{blocEdition}
\editionPage{defi-ex-3}{D03 · Défi}{Un nombre et son opposé}{Approfondissement · 30 à 60 minutes indicatives}
\begin{blocEdition}{defi-ex-3}
On cherche un nombre $a$ dont la distance à zéro est strictement comprise entre 2 et 5, et dont l’opposé est strictement négatif. Donner tous les entiers possibles, puis deux solutions décimales.\par\vspace{5mm}{\otLight\fontsize{8}{10}\selectfont\color{otMuted}Espace de recherche}\par\par\vspace{7mm}{\color{otRule!65}\hrule height .25pt}\par\vspace{7mm}{\color{otRule!65}\hrule height .25pt}\par\vspace{7mm}{\color{otRule!65}\hrule height .25pt}\par\vspace{7mm}{\color{otRule!65}\hrule height .25pt}\par\vspace{7mm}{\color{otRule!65}\hrule height .25pt}\par\vspace{7mm}{\color{otRule!65}\hrule height .25pt}
\end{blocEdition}
\editionPage{defi-ex-4}{D04 · Défi}{Remonter un carnet de bord}{Approfondissement · 30 à 60 minutes indicatives}
\begin{blocEdition}{defi-ex-4}
Un sous-marin part d’une altitude inconnue. Il descend de 35 m, remonte de 18,5 m puis descend de 7,5 m. Il termine à $-64$ m. Trouver son altitude de départ et la variation totale.\par\vspace{5mm}{\otLight\fontsize{8}{10}\selectfont\color{otMuted}Espace de recherche}\par\par\vspace{7mm}{\color{otRule!65}\hrule height .25pt}\par\vspace{7mm}{\color{otRule!65}\hrule height .25pt}\par\vspace{7mm}{\color{otRule!65}\hrule height .25pt}\par\vspace{7mm}{\color{otRule!65}\hrule height .25pt}\par\vspace{7mm}{\color{otRule!65}\hrule height .25pt}\par\vspace{7mm}{\color{otRule!65}\hrule height .25pt}
\end{blocEdition}
\editionPage{defi-ex-5}{D05 · Défi}{Une fraction entre deux autres}{Approfondissement · 30 à 60 minutes indicatives}
\begin{blocEdition}{defi-ex-5}
Trouver tous les entiers positifs $a$ tels que $3/5<a/15<4/5$. Parmi les fractions obtenues, lesquelles ont une écriture décimale finie ? Justifier par une écriture équivalente ou par la division.\par\vspace{5mm}{\otLight\fontsize{8}{10}\selectfont\color{otMuted}Espace de recherche}\par\par\vspace{7mm}{\color{otRule!65}\hrule height .25pt}\par\vspace{7mm}{\color{otRule!65}\hrule height .25pt}\par\vspace{7mm}{\color{otRule!65}\hrule height .25pt}\par\vspace{7mm}{\color{otRule!65}\hrule height .25pt}\par\vspace{7mm}{\color{otRule!65}\hrule height .25pt}\par\vspace{7mm}{\color{otRule!65}\hrule height .25pt}
\end{blocEdition}
\editionPage{defi-ex-6}{D06 · Défi}{Un ruban en trois parts}{Approfondissement · 30 à 60 minutes indicatives}
\begin{blocEdition}{defi-ex-6}
A reçoit $2/7$ d’un ruban ; B reçoit $1/3$ du ruban initial. C reçoit les 32 cm restants. Trouver la longueur initiale et les trois parts.\par\vspace{5mm}{\otLight\fontsize{8}{10}\selectfont\color{otMuted}Espace de recherche}\par\par\vspace{7mm}{\color{otRule!65}\hrule height .25pt}\par\vspace{7mm}{\color{otRule!65}\hrule height .25pt}\par\vspace{7mm}{\color{otRule!65}\hrule height .25pt}\par\vspace{7mm}{\color{otRule!65}\hrule height .25pt}\par\vspace{7mm}{\color{otRule!65}\hrule height .25pt}\par\vspace{7mm}{\color{otRule!65}\hrule height .25pt}
\end{blocEdition}
\editionPage{defi-ex-7}{D07 · Défi}{Le cadre autour d’un carré}{Approfondissement · 30 à 60 minutes indicatives}
\begin{blocEdition}{defi-ex-7}
Un carré de côté 8 cm est entouré d’une bordure extérieure de 1 cm. Retrouver l’aire de la bordure par deux méthodes, sans la compter deux fois aux coins.\label{fig-defi-7}\figureGeo{puiss-15}
\par\vspace{5mm}{\otLight\fontsize{8}{10}\selectfont\color{otMuted}Espace de recherche}\par\par\vspace{7mm}{\color{otRule!65}\hrule height .25pt}\par\vspace{7mm}{\color{otRule!65}\hrule height .25pt}\par\vspace{7mm}{\color{otRule!65}\hrule height .25pt}\par\vspace{7mm}{\color{otRule!65}\hrule height .25pt}\par\vspace{7mm}{\color{otRule!65}\hrule height .25pt}\par\vspace{7mm}{\color{otRule!65}\hrule height .25pt}
\end{blocEdition}
\editionPage{defi-ex-8}{D08 · Défi}{Retrouver l’étape d’un motif}{Approfondissement · 30 à 60 minutes indicatives}
\begin{blocEdition}{defi-ex-8}
Une rangée de carrés accolés utilise $3n+1$ allumettes à l’étape $n$. Peut-on utiliser exactement 100 allumettes ? Et 101 ? Donner une procédure arithmétique.\label{fig-defi-8}\figureGeo{lit-8}
\par\vspace{5mm}{\otLight\fontsize{8}{10}\selectfont\color{otMuted}Espace de recherche}\par\par\vspace{7mm}{\color{otRule!65}\hrule height .25pt}\par\vspace{7mm}{\color{otRule!65}\hrule height .25pt}\par\vspace{7mm}{\color{otRule!65}\hrule height .25pt}\par\vspace{7mm}{\color{otRule!65}\hrule height .25pt}\par\vspace{7mm}{\color{otRule!65}\hrule height .25pt}\par\vspace{7mm}{\color{otRule!65}\hrule height .25pt}
\end{blocEdition}
\editionPage{defi-ex-9}{D09 · Défi}{La somme de cinq consécutifs}{Approfondissement · 30 à 60 minutes indicatives}
\begin{blocEdition}{defi-ex-9}
Exprimer la somme de cinq entiers consécutifs en fonction du premier $n$, puis montrer qu’elle est cinq fois celui du milieu.\par\vspace{5mm}{\otLight\fontsize{8}{10}\selectfont\color{otMuted}Espace de recherche}\par\par\vspace{7mm}{\color{otRule!65}\hrule height .25pt}\par\vspace{7mm}{\color{otRule!65}\hrule height .25pt}\par\vspace{7mm}{\color{otRule!65}\hrule height .25pt}\par\vspace{7mm}{\color{otRule!65}\hrule height .25pt}\par\vspace{7mm}{\color{otRule!65}\hrule height .25pt}\par\vspace{7mm}{\color{otRule!65}\hrule height .25pt}
\end{blocEdition}
\editionPage{defi-ex-10}{D10 · Défi}{Un partage dessiné}{Approfondissement · 30 à 60 minutes indicatives}
\begin{blocEdition}{defi-ex-10}
Trois parts identiques et 6 euros supplémentaires donnent 45 euros. Si on ajoute ensuite 5 euros à chacune des trois parts, quel nouveau total obtient-on ? Expliquer deux méthodes.\label{fig-defi-10}\figureGeo{eq-14}
\par\vspace{5mm}{\otLight\fontsize{8}{10}\selectfont\color{otMuted}Espace de recherche}\par\par\vspace{7mm}{\color{otRule!65}\hrule height .25pt}\par\vspace{7mm}{\color{otRule!65}\hrule height .25pt}\par\vspace{7mm}{\color{otRule!65}\hrule height .25pt}\par\vspace{7mm}{\color{otRule!65}\hrule height .25pt}\par\vspace{7mm}{\color{otRule!65}\hrule height .25pt}\par\vspace{7mm}{\color{otRule!65}\hrule height .25pt}
\end{blocEdition}
\editionPage{defi-ex-11}{D11 · Défi}{Deux réductions successives}{Approfondissement · 30 à 60 minutes indicatives}
\begin{blocEdition}{defi-ex-11}
Un article coûte 100 euros. On retire d’abord 20 \% du prix, puis 10 \% du nouveau prix. Est-ce une réduction totale de 30 \% ? Justifier et déterminer la réduction totale.\par\vspace{5mm}{\otLight\fontsize{8}{10}\selectfont\color{otMuted}Espace de recherche}\par\par\vspace{7mm}{\color{otRule!65}\hrule height .25pt}\par\vspace{7mm}{\color{otRule!65}\hrule height .25pt}\par\vspace{7mm}{\color{otRule!65}\hrule height .25pt}\par\vspace{7mm}{\color{otRule!65}\hrule height .25pt}\par\vspace{7mm}{\color{otRule!65}\hrule height .25pt}\par\vspace{7mm}{\color{otRule!65}\hrule height .25pt}
\end{blocEdition}
\editionPage{defi-ex-12}{D12 · Défi}{Un modèle et son domaine}{Approfondissement · 30 à 60 minutes indicatives}
\begin{blocEdition}{defi-ex-12}
Une réserve de 24 L perd 3 L par minute. Déterminer les volumes aux minutes 3 et 7, le moment où elle est vide, puis expliquer ce que ferait un tableur si on continuait la formule après ce moment.\label{fig-defi-12}\figureGeo{rep-15}
\par\vspace{5mm}{\otLight\fontsize{8}{10}\selectfont\color{otMuted}Espace de recherche}\par\par\vspace{7mm}{\color{otRule!65}\hrule height .25pt}\par\vspace{7mm}{\color{otRule!65}\hrule height .25pt}\par\vspace{7mm}{\color{otRule!65}\hrule height .25pt}\par\vspace{7mm}{\color{otRule!65}\hrule height .25pt}\par\vspace{7mm}{\color{otRule!65}\hrule height .25pt}\par\vspace{7mm}{\color{otRule!65}\hrule height .25pt}
\end{blocEdition}
\editionPage{defi-ex-13}{D13 · Défi}{Quatre angles et une bissectrice}{Approfondissement · 30 à 60 minutes indicatives}
\begin{blocEdition}{defi-ex-13}
Autour de O, quatre angles consécutifs valent $110^\circ$, $95^\circ$, une mesure inconnue et $80^\circ$. On trace la bissectrice de l’angle inconnu. Calculer chacune de ses moitiés et expliquer pourquoi les autres angles ne changent pas.\label{fig-defi-13}\figureGeo{ang-13}
\par\vspace{5mm}{\otLight\fontsize{8}{10}\selectfont\color{otMuted}Espace de recherche}\par\par\vspace{7mm}{\color{otRule!65}\hrule height .25pt}\par\vspace{7mm}{\color{otRule!65}\hrule height .25pt}\par\vspace{7mm}{\color{otRule!65}\hrule height .25pt}\par\vspace{7mm}{\color{otRule!65}\hrule height .25pt}\par\vspace{7mm}{\color{otRule!65}\hrule height .25pt}\par\vspace{7mm}{\color{otRule!65}\hrule height .25pt}
\end{blocEdition}
\editionPage{defi-ex-14}{D14 · Défi}{Une aire en plusieurs partages}{Approfondissement · 30 à 60 minutes indicatives}
\begin{blocEdition}{defi-ex-14}
ABC a une base AB de 12 cm et une hauteur issue de C de 5 cm. M est le milieu de [AB]. N est le milieu de [AM]. Trouver les aires de ACM, ACN et NCM sans calculer de nouvelle hauteur.\label{fig-defi-14}\figureGeo{defi-triangle-partages}
\par\vspace{5mm}{\otLight\fontsize{8}{10}\selectfont\color{otMuted}Espace de recherche}\par\par\vspace{7mm}{\color{otRule!65}\hrule height .25pt}\par\vspace{7mm}{\color{otRule!65}\hrule height .25pt}\par\vspace{7mm}{\color{otRule!65}\hrule height .25pt}\par\vspace{7mm}{\color{otRule!65}\hrule height .25pt}\par\vspace{7mm}{\color{otRule!65}\hrule height .25pt}\par\vspace{7mm}{\color{otRule!65}\hrule height .25pt}
\end{blocEdition}
\editionPage{defi-ex-15}{D15 · Défi}{Un triangle et son centre}{Approfondissement · 30 à 60 minutes indicatives}
\begin{blocEdition}{defi-ex-15}
On applique un demi-tour de centre O au triangle donné. Montrer que le quadrilatère AB A′B′, pris dans cet ordre, est un parallélogramme, à condition que A, B et O ne soient pas alignés.\label{fig-defi-15}\figureGeo{sym-7}
\par\vspace{5mm}{\otLight\fontsize{8}{10}\selectfont\color{otMuted}Espace de recherche}\par\par\vspace{7mm}{\color{otRule!65}\hrule height .25pt}\par\vspace{7mm}{\color{otRule!65}\hrule height .25pt}\par\vspace{7mm}{\color{otRule!65}\hrule height .25pt}\par\vspace{7mm}{\color{otRule!65}\hrule height .25pt}\par\vspace{7mm}{\color{otRule!65}\hrule height .25pt}\par\vspace{7mm}{\color{otRule!65}\hrule height .25pt}
\end{blocEdition}
\editionPage{defi-ex-16}{D16 · Défi}{Même base, plusieurs aires}{Approfondissement · 30 à 60 minutes indicatives}
\begin{blocEdition}{defi-ex-16}
Un parallélogramme a une base de 8 cm et une hauteur de 3 cm. On déplace le côté supérieur parallèlement à la base sans changer la hauteur ni la base. L’aire change-t-elle ? Le périmètre doit-il rester identique ?\label{fig-defi-16}\figureGeo{para-4}
\par\vspace{5mm}{\otLight\fontsize{8}{10}\selectfont\color{otMuted}Espace de recherche}\par\par\vspace{7mm}{\color{otRule!65}\hrule height .25pt}\par\vspace{7mm}{\color{otRule!65}\hrule height .25pt}\par\vspace{7mm}{\color{otRule!65}\hrule height .25pt}\par\vspace{7mm}{\color{otRule!65}\hrule height .25pt}\par\vspace{7mm}{\color{otRule!65}\hrule height .25pt}\par\vspace{7mm}{\color{otRule!65}\hrule height .25pt}
\end{blocEdition}
\editionPage{defi-ex-17}{D17 · Défi}{Deux cylindres de même capacité}{Approfondissement · 30 à 60 minutes indicatives}
\begin{blocEdition}{defi-ex-17}
Un cylindre A a un rayon de 2 cm et une hauteur de 10 cm. Un cylindre B a un rayon de 4 cm. Quelle hauteur doit avoir B pour contenir le même volume ? Comparer les hauteurs.\label{fig-defi-17}\figureGeo{sol-7}
\par\vspace{5mm}{\otLight\fontsize{8}{10}\selectfont\color{otMuted}Espace de recherche}\par\par\vspace{7mm}{\color{otRule!65}\hrule height .25pt}\par\vspace{7mm}{\color{otRule!65}\hrule height .25pt}\par\vspace{7mm}{\color{otRule!65}\hrule height .25pt}\par\vspace{7mm}{\color{otRule!65}\hrule height .25pt}\par\vspace{7mm}{\color{otRule!65}\hrule height .25pt}\par\vspace{7mm}{\color{otRule!65}\hrule height .25pt}
\end{blocEdition}
\editionPage{defi-ex-18}{D18 · Défi}{Une moyenne à atteindre}{Approfondissement · 30 à 60 minutes indicatives}
\begin{blocEdition}{defi-ex-18}
Quatre notes sont 8, 12, 13 et 15. On ajoute une cinquième note. Quelle note donne une moyenne de 13 ? Quelle moyenne maximale peut-on atteindre si chaque note est au plus 20 ?\par\vspace{5mm}{\otLight\fontsize{8}{10}\selectfont\color{otMuted}Espace de recherche}\par\par\vspace{7mm}{\color{otRule!65}\hrule height .25pt}\par\vspace{7mm}{\color{otRule!65}\hrule height .25pt}\par\vspace{7mm}{\color{otRule!65}\hrule height .25pt}\par\vspace{7mm}{\color{otRule!65}\hrule height .25pt}\par\vspace{7mm}{\color{otRule!65}\hrule height .25pt}\par\vspace{7mm}{\color{otRule!65}\hrule height .25pt}
\end{blocEdition}
\editionPage{defi-ex-19}{D19 · Défi}{Deux conditions pour une urne}{Approfondissement · 30 à 60 minutes indicatives}
\begin{blocEdition}{defi-ex-19}
Une urne comporte 20 boules rouges, bleues ou vertes. La probabilité de rouge est $1/4$, celle de bleu $2/5$. Combien faut-il de boules de chaque couleur ? Vérifier que le modèle est possible.\par\vspace{5mm}{\otLight\fontsize{8}{10}\selectfont\color{otMuted}Espace de recherche}\par\par\vspace{7mm}{\color{otRule!65}\hrule height .25pt}\par\vspace{7mm}{\color{otRule!65}\hrule height .25pt}\par\vspace{7mm}{\color{otRule!65}\hrule height .25pt}\par\vspace{7mm}{\color{otRule!65}\hrule height .25pt}\par\vspace{7mm}{\color{otRule!65}\hrule height .25pt}\par\vspace{7mm}{\color{otRule!65}\hrule height .25pt}
\end{blocEdition}
\editionPage{defi-ex-20}{D20 · Défi}{Un escalier en boucle}{Approfondissement · 30 à 60 minutes indicatives}
\begin{blocEdition}{defi-ex-20}
Le script de l’escalier répète trois fois : avancer de 30, tourner à gauche de 90 degrés, avancer de 20, tourner à droite de 90 degrés. Modifier seulement le nombre de répétitions pour finir 150 pas à droite et 100 pas en haut. Donner la distance totale parcourue.\label{fig-defi-20}\figureGeo{prog-12}
\par\vspace{5mm}{\otLight\fontsize{8}{10}\selectfont\color{otMuted}Espace de recherche}\par\par\vspace{7mm}{\color{otRule!65}\hrule height .25pt}\par\vspace{7mm}{\color{otRule!65}\hrule height .25pt}\par\vspace{7mm}{\color{otRule!65}\hrule height .25pt}\par\vspace{7mm}{\color{otRule!65}\hrule height .25pt}\par\vspace{7mm}{\color{otRule!65}\hrule height .25pt}\par\vspace{7mm}{\color{otRule!65}\hrule height .25pt}
\end{blocEdition}
